\section{Panorama de las Tools integradas de Codex}

\subsection{Método para obtener el Tools Schema}

Por defecto, el Tools Schema a nivel de API de Codex no se persiste localmente. Si se necesita ver el cuerpo completo de la solicitud (incluida la definición de Tools), hay que configurar una variable de entorno específica para habilitar el registro de logs:

\begin{warningbox}{Método para volcar los logs a disco}
Al configurar la variable de entorno, Codex guarda el cuerpo completo de la solicitud a la API en la base de datos local \texttt{logs-1.sqlite}. Al consultar los registros de la base de datos se puede obtener el cuerpo completo de la solicitud, que incluye \texttt{instructions} (System Prompt), \texttt{input} (Message List) y \texttt{tools} (Tools Schema).
\end{warningbox}

Estructura del cuerpo de la solicitud (formato Response API):

\begin{itemize}
    \item \texttt{instructions}: instrucción independiente de nivel superior, equivalente al System Prompt
    \item \texttt{input}: lista completa de mensajes, que puede incluir roles como Developer, User, Assistant, Function Call Output, entre otros
    \item \texttt{tools}: lista de definiciones de herramientas
    \item Una característica propia de la Response API: se puede enlazar un Response ID anterior para encadenar el contexto
\end{itemize}

\subsection{Las cinco categorías de Tools integradas de Codex}

Al analizar el Tools Schema de Codex, se pueden identificar al menos cinco tipos de Function integradas:

\begin{importantbox}{Clasificación de las Tools integradas de Codex}
\begin{enumerate}
    \item \textbf{Ejecución de comandos de Shell}:
    \begin{itemize}
        \item \texttt{execute\_command}: ejecuta un comando de Shell
        \item \texttt{send\_input}: escribe entrada estándar en una sesión de comando existente y lee la salida estándar
    \end{itemize}
    \item \textbf{Gestión de planes}:
    \begin{itemize}
        \item \texttt{update\_plan}: actualiza el plan de ejecución (todos los Coding Agent modernos tienen Plan Mode integrado)
    \end{itemize}
    \item \textbf{Interacción con el usuario}:
    \begin{itemize}
        \item \texttt{request\_user\_input} / \texttt{ask\_user\_question}: hace una pregunta al usuario
    \end{itemize}
    \item \textbf{Operaciones sobre código}:
    \begin{itemize}
        \item \texttt{apply\_patch}: edita y ejecuta código
        \item \texttt{view\_image}: visualiza imágenes
    \end{itemize}
    \item \textbf{Red y búsqueda}:
    \begin{itemize}
        \item \texttt{websearch}: búsqueda en la red
    \end{itemize}
\end{enumerate}
\end{importantbox}

\subsection{Herramientas de Multi-Agent (subagentes)}

Codex también integra un conjunto completo de herramientas de sistema Multi-Agent:

\begin{importantbox}{Herramientas de gestión de subagentes}
\begin{itemize}
    \item \texttt{spawn\_agent}: inicia un nuevo subagente
    \item \texttt{send\_input}: envía un mensaje adicional o contexto complementario a un subagente existente
    \item \texttt{resume\_agent}: reanuda un subagente previamente cerrado o pausado
    \item \texttt{wait}: espera a que uno o varios subagentes terminen
    \item \texttt{close\_agent}: cierra un subagente y devuelve su último estado
    \item \texttt{spawn\_agents\_on\_csv}: inicia subagentes en lote según cada fila de un CSV y agrega los resultados
\end{itemize}
\end{importantbox}

\begin{knowledgebox}{Subagentes en lote a partir de CSV}
\texttt{spawn\_agents\_on\_csv} es una función relativamente reciente que permite iniciar subagentes en lote a partir de cada fila de un archivo CSV; cada subagente procesa una fila de datos y, al final, se agregan todos los resultados. Esto es útil para tareas de procesamiento por lotes y constituye un patrón práctico de orquestación de Multi-Agent.
\end{knowledgebox}

\subsection{MCP como Dynamic Tools}

Además de las Tools integradas, MCP (Model Context Protocol) ofrece un mecanismo de herramientas dinámicas:

\begin{itemize}
    \item Las Tools de MCP dependen de qué MCP Server haya configurado el usuario
    \item Cada MCP Server expone cierto número de Tools / Functions
    \item Las definiciones de Tools de MCP también se exponen al modelo a través del campo \texttt{tools} de la API
    \item MCP se configura en un archivo de configuración, pero el Tools Schema que expone no aparece en el archivo de Session
\end{itemize}

\subsection{Resumen de la sección}

Codex integra cinco categorías de Tools (ejecución de Shell, gestión de planes, interacción con el usuario, operaciones sobre código, búsqueda en la red), además de un conjunto completo de herramientas de gestión de subagentes de Multi-Agent. MCP ofrece capacidad de extensión de herramientas dinámicas. Todas las Tools se exponen al modelo a través del campo tools de la API.
